\documentclass{exam}

\newif\ifA
\newif\ifkey
\newif\ifprintselected   % required by exam_config's \ansbox

%%%%%%%%%%%%%%%%%%%%%%%%%%%%%%%%%%%%%%%%%%%%%%%%%%%%%%
%%% SET THESE IF VARIABLES
%%%%%%%%%%%%%%%%%%%%%%%%%%%%%%%%%%%%%%%%%%%%%%%%%%%%%%
% Lab name -- labs are named by the lesson they belong to, never by a
% sequential lab number (Neil, 2026-10-12): this is the L14 Lab.
\def\labLesson{14}

% is this the key or not?
\ifdefined\iskey
	\ifnum\iskey=1
		\keytrue
	\else
		\keyfalse
	\fi
\else
	%% Pick one manually
	\keyfalse % if this is NOT the key
%	\keytrue     % if this IS the key
\fi

% set up the header and footer - change for every term
\ifdefined\thisTermIn
	\def\thisterm{\thisTermIn}
\else
	\def\thisterm{Spring 2027}
\fi
%
\def\thisclass{ECE 448}

%%%%%%%%%%%%%%%%%%%%%%%%%%%%%%%%%%%%%%%%%%%%%%%%%%%%%%
%%% END SET THESE IF VARIABLES
%%%%%%%%%%%%%%%%%%%%%%%%%%%%%%%%%%%%%%%%%%%%%%%%%%%%%%
% packages and shortcuts
\input{myPackages_gr}
\usepackage[hidelinks, linkcolor=red]{hyperref}
\input{myShortcuts}
\setlength{\parindent}{0pt}
\setlength{\parskip}{6pt}

%% exam class customization
\input{exam_config}
% ECE 448 document fonts.
%
% Body text is Barlow (SIL OFL; the faces and the license are in
% latex/fonts/Barlow/). Barlow is not in TeX Live, so it is loaded as an
% OpenType face through fontspec -- which means the document must be built
% with lualatex, not pdflatex. build_practice.sh and build_lab.sh both do.
%
% Compiling with pdflatex anyway is not an error: the guard below falls back
% to the LaTeX default so the source still compiles under any engine and can
% be syntax-checked. The resulting PDF just is not the one we ship.
%
% Paths are relative to latex/, which is where both build scripts run
% (each does cd "$(dirname "$0")"). Loading by path rather than by installed
% family name means no system font install and no fontconfig, so a fresh
% clone builds a byte-comparable PDF anywhere.
%
% MATH IS DELIBERATELY LEFT IN COMPUTER MODERN. Barlow has no math companion,
% and CM math under sans body text is already the house pairing: the decks set
% words in Source Sans Pro and let MathJax render symbols in its CM-derived
% face. Moving math to a sans as well would mean unicode-math plus a sans math
% font (Fira Math being the plausible one) and re-checking every equation in
% the course for symbol coverage. That is a separate decision, not a side
% effect of this one.

\usepackage{iftex}
\ifLuaTeX
	\usepackage{fontspec}
	\defaultfontfeatures{Ligatures=TeX}
	\setmainfont{Barlow}[
		Path           = fonts/Barlow/,
		Extension      = .ttf,
		UprightFont    = *-Regular,
		ItalicFont     = *-Italic,
		BoldFont       = *-Bold,
		BoldItalicFont = *-BoldItalic,
	]
	\setsansfont{Barlow}[
		Path           = fonts/Barlow/,
		Extension      = .ttf,
		UprightFont    = *-Regular,
		ItalicFont     = *-Italic,
		BoldFont       = *-SemiBold,
		BoldItalicFont = *-BoldItalic,
	]
	% microtype earns its place here specifically because of Barlow. Its
	% glyphs are narrower than Computer Modern's, so justified prose sets
	% tighter and TeX has to squeeze interword space to fit; two paragraphs
	% in the L02/L03 sets went overfull on the first Barlow build. Protrusion
	% and font expansion absorb that -- both are lualatex features, so this
	% belongs inside the guard with everything else.
	\usepackage{microtype}

	% Mono stays Latin Modern: Barlow ships no monospaced face, and verbatim
	% and code need one that actually is fixed-pitch.
\fi
   % Barlow body text (needs lualatex)
\printselectedfalse      % full worked solutions in the key, not answers-only
\CorrectChoiceEmphasis{\color{red}}
\SolutionEmphasis{\color{red}}
\renewcommand\questionlabel{\textbf{\thequestion.}}
\renewcommand{\thepartno}{\alph{partno}}
\renewcommand{\partlabel}{(\thepartno)}

% Fillable table cell: blank in the student copy, the expected value in red on
% the KEY. Fixed-width centered columns keep the blank cells writable.
\newcolumntype{K}[1]{>{\centering\arraybackslash}p{#1}}
\newcommand{\kv}[1]{\ifkey{\color{red}#1}\fi\rule[-0.10in]{0pt}{0.34in}}

% Questions 1 to 3 (on the test file), Question 4 (the lab capture), and Question 7(b) take their answers from
% L14_capture_key.tex, which lives in ECE448-faculty and is never committed
% here. Without it, even a KEY build prints placeholders.
\newcommand{\capfk}{(faculty key)}
\newcommand{\qOneUpper}{\capfk}\newcommand{\qOneLower}{\capfk}\newcommand{\qOneDiff}{\capfk}
\newcommand{\qOneSol}{\capfk}\newcommand{\qOneShot}{\capfk}
\newcommand{\qTwoCell}{\capfk}\newcommand{\qTwoSol}{\capfk}
\newcommand{\capEdgeLo}{\capfk}\newcommand{\capEdgeLi}{\capfk}
\newcommand{\capEdgeUi}{\capfk}\newcommand{\capEdgeUo}{\capfk}
\newcommand{\capUpperHybrid}{\capfk}\newcommand{\capLowerHybrid}{\capfk}
\newcommand{\capInnerStop}{\capfk}
\newcommand{\capUpperInner}{\capfk}\newcommand{\capLowerInner}{\capfk}
\newcommand{\capDiff}{\capfk}
\newcommand{\capSolEdges}{\capfk}\newcommand{\capSolMeters}{\capfk}
\newcommand{\capSolTotal}{\capfk}\newcommand{\capSolNeighbor}{\capfk}
\newcommand{\capSolWriteup}{\capfk}
\IfFileExists{L14_capture_key.tex}{\input{L14_capture_key}}{}

% print the answers if this is the key
\ifkey
	\printanswers % if this IS the key
\else
	\noprintanswers  %if this is NOT the key
\fi

% print the header and footer
% need to print key on top if this is the key
\ifkey
	\firstpageheader{}{\color{red}{ KEY}}{}
	\runningheader{\thisclass\hspace{1pt} -- \thisterm }{L\labLesson\ Lab\hspace{1pt}\color{red}{ KEY}}{\thepage\hspace{1pt} of\hspace{1pt} \numpages}
	\firstpagefooter{}{}{}
	\runningfooter{}{}{}
\else
	\firstpageheader{}{}{}
	\runningheader{\thisclass \hspace{1pt} -- \thisterm }{L\labLesson\ Lab\hspace{1pt}}{\thepage\hspace{1pt} of\hspace{1pt} \numpages}
	\firstpagefooter{}{}{}
	\runningfooter{}{}{}
\fi

\runningheadrule




\begin{document}
\pagestyle{headandfoot}   % exam class

%%%%%%%%%%%%%%%%%%%%%%%%%%%%%%%%%%%%%%%%%%%%%%%%%%%%%%
% title block
%%%%%%%%%%%%%%%%%%%%%%%%%%%%%%%%%%%%%%%%%%%%%%%%%%%%%%

\begin{center}
USAF ACADEMY \\
DEPARTMENT OF ELECTRICAL AND COMPUTER ENGINEERING \\
LESSON \labLesson\ LAB - HD RADIO \\
\textbf{Lab Packet} \\
\thisterm \\[4mm]
\end{center}

\vspace{-4pt}
\textbf{Name:} \rule[-2pt]{2.6in}{0.4pt} \hfill
\textbf{Section:} \rule[-2pt]{0.9in}{0.4pt} \hfill
\textbf{Date:} \rule[-2pt]{1.1in}{0.4pt}

\vspace{4pt}

\fullwidth{\textbf{Learning Objective 6.3: } I can describe how HD Radio works and
demodulate an HD Radio signal.}

\vspace{2pt}

This packet is the L\labLesson\ Lab, in three parts, and you turn it in once, complete,
before Lesson 15. Do \textbf{Part 1} on your own before the lab period: run the sideband
finder, change one test-file setting, and slide Inner subcarrier to see what it does. \textbf{Part 2} is your own
measurements, recorded in class on your own laptop with your RTL-SDR. Do \textbf{Part 3}
on your own after the lab: explain what you measured and answer the
analysis questions.

\textbf{Grading.} One submission, graded against Objective 6.3 on the course's
1 / 0 scale, with revision allowed under course policy. Paste screenshots where a
question asks for one.

\textbf{Transmitting.} Nothing in this lab transmits. The RTL-SDR receives only, and both
captures came from GNU Radio flowgraphs, not from an antenna.

The flowgraphs are on the L14 lesson page, \texttt{L14-hd-test-tx.grc} and
\texttt{L14-hd-find.grc}, and the capture \texttt{L14-lab-capture.cu8}, which your instructor posts on
the L14 Lab page before the lab. Both captures are simulated: unsigned 8-bit I and Q at 1,488,375 samples per second,
the format the RTL-SDR and nrsc5 record. Every number below uses the hybrid values from
the lesson: subcarriers 363.4 Hz apart, one FFT bin at \texttt{fft\_len} 4096, sidebands from
subcarrier 356 to 546 on each side, partitions of 19 subcarriers, and an extended hybrid
station adding one, two, or four partitions inside subcarrier 356 on each side.

\begin{questions}

\fullwidth{\large\textbf{Part 1: Pre-Lab} \normalsize(individual, before the lab period)}

%%%%%%%%%%%%%%%%%%%%%%%%%%%%%%%%%%%%%%%%%%%%%%%%%%%
\question \textbf{Change a setting and read the finder.} Put both L14 flowgraphs in one
folder. Open \texttt{L14-hd-test-tx.grc}, change the Variable \texttt{upper\_dbc} from $-23$
to $-18$, and run it once to write \texttt{L14-hd-test.cu8}. Then open and run
\texttt{L14-hd-find.grc}.

\begin{parts}

	\begin{minipage}[t]{0.58\linewidth}
		\part After about 5 s, record the three meter readings.
		\begin{solution}[0.6in]
		\qOneSol
		\end{solution}
	\end{minipage}\hfill%
	\begin{minipage}[t]{0.37\linewidth}
		\raggedleft \vspace{12pt}
		upper \ansbox[1.1in]{\qOneUpper}

		\vspace{4pt}
		lower \ansbox[1.1in]{\qOneLower}

		\vspace{4pt}
		upper $-$ lower \ansbox[1.1in]{\qOneDiff}
	\end{minipage}

	\part Paste a screenshot of the running flowgraph's window. What changed on the frequency
	sink compared with the lesson's figure?
		\begin{solution}[1.6in]
		\qOneShot
		\end{solution}

\end{parts}

%%%%%%%%%%%%%%%%%%%%%%%%%%%%%%%%%%%%%%%%%%%%%%%%%%%
\newpage
\question \textbf{The Inner subcarrier slider.} Keep the finder running on the same file.
Slide \textbf{Inner subcarrier} from 356 down to 280, one step at a time, and record the
upper sideband reading at each step.

	\begin{center}
	\renewcommand{\arraystretch}{1.25}
	\begin{tabular}{| l | K{0.8in} | K{0.8in} | K{0.8in} | K{0.8in} | K{0.8in} |}
	\hline
	Inner subcarrier & 356 & 337 & 318 & 299 & 280 \\\hline
	Upper (dBc) & \kv{\qTwoCell} & \kv{\qTwoCell} & \kv{\qTwoCell} & \kv{\qTwoCell} & \kv{\qTwoCell} \\\hline
	\end{tabular}
	\end{center}

	\begin{solution}[0.8in]
	\qTwoSol
	\end{solution}

%%%%%%%%%%%%%%%%%%%%%%%%%%%%%%%%%%%%%%%%%%%%%%%%%%%
\newpage
\fullwidth{\large\textbf{Part 2: Measurements} \normalsize(in class)}

\question \textbf{Build the sideband meter.} Start GRC with an empty canvas. Do not open
\texttt{L14-hd-find.grc}: build the meter yourself as the lab page lists.

\begin{parts}

	\part Paste a screenshot of your flowgraph on the GRC canvas. Every connection must be
	black, not red.
		\begin{solution}[1.6in]
		The flowgraph matches the lab page: the byte File Source, Throttle, UChar To Float,
		Add Const, Multiply Const, Deinterleave, and Float To Complex; a frequency sink; three
		Frequency Xlating FFT Filters, each followed by Complex to Mag\^{}2, Moving Average,
		Keep 1 in N, and Log10; three Subtract blocks; and a Number Sink with three inputs.
		\end{solution}

	\begin{minipage}[t]{0.58\linewidth}
		\part Run it on the file you wrote in Question 1 and record the three readings. Paste
		a screenshot of the running window.
		\begin{solution}[1.2in]
		Within 0.1 dB of the finder's readings in Question 1(a), because the meter uses the
		same filters and averaging.
		\end{solution}
	\end{minipage}\hfill%
	\begin{minipage}[t]{0.37\linewidth}
		\raggedleft \vspace{12pt}
		upper \ansbox[1.1in]{\qOneUpper}

		\vspace{4pt}
		lower \ansbox[1.1in]{\qOneLower}

		\vspace{4pt}
		upper $-$ lower \ansbox[1.1in]{\qOneDiff}
	\end{minipage}

\end{parts}

%%%%%%%%%%%%%%%%%%%%%%%%%%%%%%%%%%%%%%%%%%%%%%%%%%%
\newpage
\question \textbf{Identify the lab capture.} Run your meter on
\texttt{L14-lab-capture.cu8}, as the lab page lists. Treat the center as 100.3 MHz.

\begin{parts}

	\part Zoom in on each sideband's edges with an FFT of 8192 points and read them. Are the
	inner edges at 129.4 kHz from the center, or closer to the host?

	\begin{center}
	\renewcommand{\arraystretch}{1.25}
	\begin{tabular}{| l | K{1.0in} | K{1.0in} | K{1.0in} | K{1.0in} |}
	\hline
	 & Lower outer & Lower inner & Upper inner & Upper outer \\\hline
	Edge (MHz) & \kv{\capEdgeLo} & \kv{\capEdgeLi} & \kv{\capEdgeUi} & \kv{\capEdgeUo} \\\hline
	\end{tabular}
	\end{center}

		\begin{solution}[0.6in]
		\capSolEdges
		\end{solution}

	\part Record both meter readings with Inner subcarrier at 356. If the inner edge is not at
	129.4 kHz, slide Inner subcarrier down until the readings stop rising, and record the
	setting and both readings there.

	\begin{center}
	\renewcommand{\arraystretch}{1.25}
	\begin{tabular}{| l | K{1.4in} | K{1.4in} |}
	\hline
	Inner subcarrier & Upper (dBc) & Lower (dBc) \\\hline
	356 & \kv{\capUpperHybrid} & \kv{\capLowerHybrid} \\\hline
	\kv{\capInnerStop} & \kv{\capUpperInner} & \kv{\capLowerInner} \\\hline
	\end{tabular}
	\end{center}

		\begin{solution}[0.5in]
		\capSolMeters
		\end{solution}

	\begin{minipage}[t]{0.58\linewidth}
		\part Do the two sidebands read the same? If not, which is stronger, and by how many
		dB on the Upper minus lower meter?
		\begin{solution}[1.0in]
		\capSolTotal
		\end{solution}
	\end{minipage}\hfill%
	\begin{minipage}[t]{0.37\linewidth}
		\raggedleft \vspace{12pt}
		difference \ansbox[1.1in]{\capDiff}
	\end{minipage}

	\part Find the other station in the capture, give its offset from the center, and say
	whether it carries HD Radio. How can you tell?
		\begin{solution}[0.6in]
		\capSolNeighbor
		\end{solution}

\end{parts}

%%%%%%%%%%%%%%%%%%%%%%%%%%%%%%%%%%%%%%%%%%%%%%%%%%%
\newpage
\question \textbf{A live station.} Add the RTL-SDR Source to your meter and find a local
station with HD Radio, as the lab page lists.

\begin{parts}

	\part Paste a screenshot with the shelves visible, and record the station, the RF
	gain, and the height of the shelves above the floor.
		\begin{solution}[1.6in]
		Answers vary. The shelves should look like the test file's, with a rougher host whose
		width changes with the audio on the waterfall.
		\end{solution}

	\part Record the inner and outer edges in MHz, the Inner subcarrier setting where the
	readings stop rising, and the configuration.
		\begin{solution}[0.8in]
		Answers vary. Most stations are hybrid: edges about 129.4 and 198.4 kHz from the
		center, and readings that do not rise below 356. Edges closer to the host, with
		readings that rise as the slider moves down, mean extended hybrid.
		\end{solution}

	\begin{minipage}[t]{0.58\linewidth}
		\part Record both meter readings. Then lower the gain in steps and record what happens
		to the shelves and the readings as the floor rises toward them.
		\begin{solution}[0.9in]
		Answers vary; a station at the original level reads about $-23$ dBc per sideband,
		and a stronger one reads higher. As the floor rises toward the shelves, the noise in
		each band adds to the reading, so the meter reads high, and the shelves' edges blur
		into the floor.
		\end{solution}
	\end{minipage}\hfill%
	\begin{minipage}[t]{0.37\linewidth}
		\raggedleft \vspace{12pt}
		upper \ansbox[1.1in]{}

		\vspace{4pt}
		lower \ansbox[1.1in]{}
	\end{minipage}

\end{parts}

%%%%%%%%%%%%%%%%%%%%%%%%%%%%%%%%%%%%%%%%%%%%%%%%%%%
\newpage
\question \textbf{Decode it.} Decode the same station with nrsc5, as the lab page lists.
These steps are documented usage from the nrsc5 README; we have not yet tested them in
this course.

\begin{parts}

	\part Record the station name, the title and artist of what is playing, and the
	programs the station carries.
		\begin{solution}[1.0in]
		Answers vary. nrsc5 prints the station name and the title and artist as they arrive,
		and lists the station's audio programs; HD1 is program 0.
		\end{solution}

	\part Lower the gain until the audio stops. At what gain does it stop? Then stop nrsc5,
	run your meter at that gain, and record how high the shelves stand above the floor.
		\begin{solution}[1.0in]
		Answers vary. The digital audio stops abruptly, not gradually as analog FM does, while
		the shelves still stand a little above the floor.
		\end{solution}

	\begin{minipage}[t]{0.58\linewidth}
		\part Record 10 s with \texttt{nrsc5 -w}, open the recording in your meter, and
		record both readings. Do they match Question 5(c)?
		\begin{solution}[0.9in]
		Answers vary. They should be close if the gain is the same, because the file holds the
		same samples the RTL-SDR delivers live.
		\end{solution}
	\end{minipage}\hfill%
	\begin{minipage}[t]{0.37\linewidth}
		\raggedleft \vspace{12pt}
		upper \ansbox[1.1in]{}

		\vspace{4pt}
		lower \ansbox[1.1in]{}
	\end{minipage}

\end{parts}

%%%%%%%%%%%%%%%%%%%%%%%%%%%%%%%%%%%%%%%%%%%%%%%%%%%
\newpage
\fullwidth{\large\textbf{Part 3: Writeup} \normalsize(individual, after the lab)}

\question \textbf{Analysis.}

\begin{parts}

	\part Compare the live station with the lesson's test file: where its shelves sit,
	how high they stand, and whether the two sides match. What does each difference tell
	you about the station's HD Radio signal?
		\begin{solution}[1.5in]
		Answers vary. Edges at 129.4 and 198.4 kHz mean hybrid; an inner edge closer to the
		host means extended hybrid. Shelves that read higher in dBc than the test file's mean
		the station runs its digital signal stronger than the test file. Two sides that read
		differently mean the station runs one sideband weaker than the other.
		\end{solution}

	\part Explain how the lab capture's waterfall, edges, and meter readings identified its
	configuration.
		\begin{solution}[1.5in]
		\capSolWriteup
		\end{solution}

	\part nrsc5 cannot decode \texttt{L14-hd-test.cu8}, although its sidebands have the right
	subcarriers, spacing, and guard. Name two things nrsc5 needs from a signal that the test
	file does not provide.
		\begin{solution}[1.5in]
		Any two: reference subcarriers that carry NRSC-5's known synchronization and control
		bits, which nrsc5 uses to find the timing, frequency, and phase; data subcarriers
		that carry real NRSC-5 frames, scrambled, coded, and interleaved as the standard
		specifies, rather than random QPSK; and compressed audio and station data inside
		those frames.
		\end{solution}

\end{parts}

\end{questions}

\vspace{1.0em}
\noindent\textbf{Documentation:} \rule[-2pt]{0.6\linewidth}{0.4pt}

\end{document}
